\documentclass[aspectratio=169]{ctexbeamer}
\usepackage[T1]{fontenc}
\usepackage{latexsym,amsmath,xcolor,multicol,booktabs,calligra}
\usepackage{graphicx,listings,stackengine}
\usepackage{hyperref}

\usepackage{PekingU}

% 去掉每节/子节前面的自动目录页
\AtBeginSection[]{}
\AtBeginSubsection[]{}

% ===== 每节课改这里 =====
\author{Cap1taL}
\title{简单网络流、二分图、与图论杂题}
\subtitle{我求求你别再追查牛战士了，行吗？我打了32把，2700的AK，我用了多少把，经济吃没了，队友被我吃垮了，现在好不容易有了便宜枪，你们非说它是瘤子，那枪瘤不瘤我们能不知道吗？那枪才1400一把，枪贩子根本不挣钱，谁家能不遇上个eco局，你就能保证你这一辈子都有钱起AK吗？你把他骂走了，我们都得等死。我不想死，我想活着，行吗？}
\date{2026年8月}
% ========================

% 代码高亮设置
\definecolor{deepblue}{rgb}{0,0,0.5}
\definecolor{deepred}{rgb}{0.6,0,0}
\definecolor{deepgreen}{rgb}{0,0.5,0}
\definecolor{halfgray}{gray}{0.55}

\lstset{
    basicstyle=\ttfamily\small,
    keywordstyle=\bfseries\color{deepblue},
    emphstyle=\ttfamily\color{deepred},
    stringstyle=\color{deepgreen},
    numbers=left,
    numberstyle=\small\color{halfgray},
    rulesepcolor=\color{red!20!green!20!blue!20},
    frame=shadowbox,
}

\begin{document}

\kaishu

% 封面
\begin{frame}
    \titlepage
\end{frame}

% 目录
% \begin{frame}{目录}
%     \tableofcontents[sectionstyle=show,subsectionstyle=show/shaded/hide,subsubsectionstyle=show/shaded/hide]
% \end{frame}

% ===== 从这里开始写你的内容 =====

\section{网络流}

\subsection{基本算法}
\begin{frame}{最大流-基本概念}
    \begin{itemize}
        \item 最大流最小割定理
        \item 残量网络
        \item 增广路
        \item 反向边与退流的作用
        \item 不断增广得到最大流
        \item 朴素增广的复杂度可以达到 $O(mf)$
    \end{itemize}
\end{frame}

\begin{frame}{最大流-Dinic}
    \begin{itemize}
        \item 给残量网络 bfs 分层
        \item 从源点开始在层间 “推流”，寻找阻塞流
        \item 每个节点尝试所有出边推流，推流大小为剩余入流量与 cap 的最小值
        \item 当前弧优化
        \item Dinic 的最坏复杂度是 $O(n^2m)$
        \item 在单位网络上为 $m\min(m^{\frac{1}{2}},n^{\frac{2}{3}})$ 
    \end{itemize}
\end{frame}

\begin{frame}{最大流-ISAP}
    \begin{itemize}
        \item Dinic 每次都要重新 bfs，有没有更厉害的算法
        \item 有的兄弟，有的
        \item 与 Dinic 不同，ISAP 采取一次 bfs 的策略。在算法初始时从 T 到 S 分层，并维护各个层的点数
        \item 当一个点增广结束时，意味着它无法在当前层获得更多增广，于是把它的层数减一，同步维护各层的点数
        \item 推流仍然只能跨一层进行。所以显然当某一层的点数为 $ 0 $ 时，我们无法在残量网络上进行任何增广，从而得到最大流
        \item ISAP 的渐进最坏时间复杂度与 Dinic 相同，但实际使用一般会快一半以上
    \end{itemize}
\end{frame}

\begin{frame}{最小费用最大流-SSP}
    \begin{itemize}
        \item 适用于无负环的情况
        \item 每次增广的时候改为寻找费用最小的增广路
        \item 只需要把 Dinic 的 bfs 换成 SPFA
        \item 增广时从必须跨层增广变为必须沿着最短路 DAG 增广
        \item 复杂度为 $O(nmf)$
    \end{itemize}
\end{frame}

% \subsection{P2053 [SCOI2007] 修车}

% \begin{frame}{P2053 [SCOI2007] 修车}
%     $M$ 位技术人员，$N$ 位车主。第 $j$ 位技术人员维修第 $i$ 辆车需要时间 $T_{i,j}$。
%     需要安排每位技术人员所维修的车辆及顺序，使所有顾客的平均等待时间最小。
%     其中，顾客的等待时间指从送车到维修完毕所用的总时间。

%     \medskip
%     $2\le M\le 9$，$1\le N\le 60$，$1\le T_{i,j}\le 10^3$。
% \end{frame}

% \begin{frame}{Solution}
%     \begin{itemize}
%         \item 
%     \end{itemize}
% \end{frame}

\section{二分图}

\subsection{常用算法}
\begin{frame}{二分图相关概念}
    \begin{itemize}
        \item 最小点覆盖等于最大匹配
        \item 最大独立集等于最小点覆盖
        \item 完美匹配
        \item 最大匹配
    \end{itemize}
\end{frame}

\begin{frame}{二分图最大匹配}
    \begin{itemize}
        \item 网络流算法运行更快，且 Dinic 并不难写
        \item 建立源点，连向左部
        \item 建立汇点，连向右部
        \item 单位网络，复杂度为 $O(\sqrt{n}m)$
    \end{itemize}
\end{frame}

\begin{frame}{Hall 定理}
    \begin{itemize}
        \item Hall 定理给出了二分图完美匹配存在的充要条件
        \item 不妨设左部点更少。若对于其任意一个子集，邻域大小都大于集合大小，则完美匹配存在
        \item 必要性显然正确，若存在一个子集不满足上述条件，则这个子集内就会出现失配
        \item 考虑利用数学归纳证明充分性。设 $N(S)$ 表示子集 $S$ 的邻域
        \begin{itemize}
            \item $n=1$ 的边界显然成立
            \item 如果存在一个真子集 $|S|=N(S)$，则我们可以删除这个子集。假设删除后出现一个非法子集 $T$，则我们可以在一开始选取 $S$ 与 $T$ 的交集，同样会产生矛盾。
            \item 如果所有子集都满足 $|S|> N(S)$，则我们可以随便删掉一对点，仍然满足 $|S|\ge N(S)$，从而归纳到了子问题
        \end{itemize}
        \item 综上 Hall 定理成立
    \end{itemize}
\end{frame}


\subsection{Codeforces 1634E Fair Share}

\begin{frame}{Codeforces 1634E Fair Share}
    给定 $m$ 个数组，第 $i$ 个数组包含 $n_i$ 个整数，且 $n_i$ 为偶数。有两个初始为空的可重集 $L, R$。要求将每个数组中的数平均分成两组，分别放入 $L$ 和 $R$，使得最终 $L=R$。判断是否存在这样的分配方案。

    \medskip
    $1\le m\le 10^5$，$2\le n_i$ 且 $2\mid n_i$，$\sum n_i\le 2\times 10^5$，$1\le a_i\le 10^9$。
\end{frame}

\begin{frame}{Solution}
    \begin{itemize}
        \item 首先判掉数有奇数个的情况，显然无解
        \item 操作非常自由，看起来非常可以乱搞，考虑加强限制
        \item 我们钦定一个序列里，每两个数不能位于同一个集合，在他们之间连边
        \item 再对于每个数，每两个连边
        \item 可以发现图总是由偶环构成的。假如存在奇环，必然存在一个点，有两条位置限制连边或值限制连边，而这是不可能的
        \item 因此这种情况下总是有解，二分图染色即可
    \end{itemize}
\end{frame}

\subsection{Codeforces 1519F Chests and Keys}

\begin{frame}{Codeforces 1519F Chests and Keys}
    有 $n$ 个宝箱，第 $i$ 个宝箱中有 $a_i$ 枚金币；有 $m$ 把钥匙，第 $j$ 把钥匙售价为 $b_j$。
    Alice 可以在第 $i$ 个宝箱上安装若干种锁，其中第 $j$ 种锁只能用第 $j$ 把钥匙打开，安装花费为 $c_{i,j}$。
    Bob 可以购买任意数量的钥匙，然后打开他能打开的所有宝箱。若 Bob 的收益（打开宝箱的金币总数减去购买钥匙的费用）为正，则 Bob 获胜；否则 Alice 获胜。
    Alice 希望加锁后，无论 Bob 如何购买钥匙，Bob 都无法获得正收益。求满足条件的最少加锁花费；若不存在方案，输出 $-1$。
    
    \medskip
    $1\le n,m\le 6$，$1\le a_i,b_j\le 4$，$1\le c_{i,j}\le 10^7$。
\end{frame}

\begin{frame}{Solution}
    \begin{itemize}
        \item 设 $L_i$ 表示 $i$ 号宝箱上锁的集合，则对于任何一个开宝箱集合 $X$，我们都需要满足：
        $$
            \sum_{i\in L_{i\in X}} b_i \ge \sum_{i\in X} a_i
        $$
        \item 类似 Hall 定理，因此我们拆点，把每个箱子拆成 $a_i$ 个点，每个锁拆成 $b_i$ 个点，若上面有锁就全部连边。我们要求它存在完美匹配
        \item 直接设 $f_{i,S}$ 表示宝箱侧前 $i$ 个点，每个锁还有几个点没匹配的集合 $S$ 的答案
        \item 转移只需要枚举当前宝箱点匹配什么点即可，复杂度 $O(n\times 5^{2n})$，DP 的常数很小
    \end{itemize}
\end{frame}

% ===== 结束页 =====

\begin{frame}
    \begin{center}
        {\Huge\calligra Thanks!}
    \end{center}
\end{frame}

\end{document}
